\chapter{Capstone 试教材料与课程运行包}

\section{案例导读：试教是在真实课堂前做一次压力测试}

一套材料在作者电脑上跑通，并不等于它能在课堂中被不同背景的学生读懂、运行、提交和接受反馈。课程材料的 bug 经常不在代码里，而在入口不清、时间估计错误、rubric 解释不足或隐私边界没有说清。

本章把试教看成课程运行前的最小实验。它要求我们像测试代码一样测试教学材料：学生是否知道下一步、助教是否能稳定评分、notebook 是否能复跑、反馈是否能进入修订闭环。

\section{案例目标}

第 39 章到第 42 章已经把 Capstone 的选题、执行、最终交付和评分 rubric 连成一条线。本章处理另一个常被低估的问题：\textbf{这些材料真的放进一门课以后，学生、助教和教师能不能顺利运行。}

本章的目标有四点：

\begin{enumerate}
    \item 理解试教不是简单地“找几名学生跑 notebook”，而是一次课程运行压力测试；
    \item 学会检查学生 handout、notebook runtime、rubric calibration、TA feedback 和 privacy boundary；
    \item 学会区分“可以正式运行”“需要修订材料”“需要校准评分”和“试教前阻塞”四类出口；
    \item 学会把试教反馈整理成下一轮课程修订任务，而不是只收集零散意见。
\end{enumerate}

\section{课程材料真正的 bug 往往不在代码里}

一份 notebook 能在作者电脑上运行，并不意味着它已经可以进入课堂。真实课堂里会出现更多变量：

\begin{itemize}
    \item 学生是否知道第一步应该做什么；
    \item handout 是否把交付物、截止时间和评分标准说清楚；
    \item notebook 是否能在课程环境里从头运行；
    \item 助教是否能用同一套 rubric 给出一致反馈；
    \item 学生材料里是否包含不能上传、不能共享或不能交给外部 AI 服务的内容；
    \item 反馈收集之后，课程团队是否知道下一轮该改哪里。
\end{itemize}

所以，本章的核心立场是：\textbf{试教不是一次展示，而是一次课程系统调试。}

\section{一个极小的 trial-teaching feedback 数据集}

配套 notebook 使用 \nolinkurl{capstone_trial_teaching_feedback_demo.csv}。这是一组完全本地、教学用途的\textbf{试教运行记录卡片}。每一行代表一次 Capstone 课堂活动或试教环节，包含：

\begin{itemize}
    \item \texttt{trial\_id}；
    \item \texttt{session\_theme}；
    \item \texttt{student\_readiness\_status}；
    \item \texttt{handout\_clarity\_status}；
    \item \texttt{notebook\_runtime\_status}；
    \item \texttt{rubric\_calibration\_status}；
    \item \texttt{ta\_feedback\_status}；
    \item \texttt{privacy\_boundary\_status}；
    \item \texttt{revision\_load\_score}；
    \item \texttt{reference\_route}。
\end{itemize}

这里的 \texttt{reference\_route} 分成四类：

\begin{itemize}
    \item \texttt{ready\_for\_rollout}：可以进入正式课堂运行；
    \item \texttt{revise\_materials}：学生 handout、TA feedback 或材料组织需要修改；
    \item \texttt{calibrate\_rubric}：评分标准、样例反馈或助教尺度还需要校准；
    \item \texttt{blocked\_before\_trial}：notebook 无法运行或隐私边界不清，试教前必须先停下。
\end{itemize}

这个数据集故意放入了几类常见试教问题：材料看起来简单但学生读不懂、notebook 在作者机器外变慢或失败、助教反馈尺度不一致、以及 AI 使用边界没有写清楚。

\section{先看一个常见 baseline：只按 revision load 排序}

课程团队很容易优先处理“看起来最省事”的环节。一个 naive baseline 是按修订负担分数从低到高排序，认为修订负担最低的环节最适合正式 rollout。

\begin{examplebox}[只按 revision load 取前三个试教环节]
\begin{lstlisting}[style=pythonstyle]
top3 = sorted(
    rows,
    key=lambda row: row["revision_load_score"],
)[:3]
\end{lstlisting}
\end{examplebox}

这个 baseline 的问题是：修订负担低不代表材料可以运行。有些环节只需要补 TA feedback 格式，看起来负担不高；但如果助教不能统一反馈，正式课程里仍然会出现评分漂移。

在本章教学数据上，只按 revision load 取前三个环节时，可正式运行的精度是 0.67。表~\ref{tab:ch43_revision_load_vs_trial_workflow} 展示了结构化 trial workflow 如何先检查 privacy boundary 和 notebook runtime，再处理 rubric calibration 与材料清晰度：

\begin{table}[htbp]
\centering
\small
\setlength{\tabcolsep}{4pt}
\begin{tabular}{@{}p{0.27\textwidth}>{\centering\arraybackslash}p{0.18\textwidth}>{\centering\arraybackslash}p{0.18\textwidth}p{0.29\textwidth}@{}}
\toprule
方法 & 可正式运行精度 & 阻塞项混入率 & 教学解读 \\
\midrule
只按 revision load 取前三项 & 0.67 & 0.00 & 低修订负担不等于助教反馈和材料清晰度已经过关 \\
结构化 trial workflow & 1.00 & 0.00 & 先排除阻塞，再校准 rubric 和材料 \\
\bottomrule
\end{tabular}
\caption{本章 notebook 中两个最小试教流程的对照。重点不是让试教流程更官僚，而是避免把课堂运行风险藏到正式开课后。}
\label{tab:ch43_revision_load_vs_trial_workflow}
\end{table}

\section{一个可执行的 trial workflow}

本章 notebook 的 routing 逻辑可以分成四层：

\begin{enumerate}
    \item \textbf{边界 gate}：隐私、共享和外部 AI 服务边界不清楚时，不进入试教；
    \item \textbf{运行 gate}：notebook 在课程环境里不能从头运行时，不进入试教；
    \item \textbf{评分校准 gate}：rubric 未校准时，先做助教 norming session；
    \item \textbf{材料修订 gate}：handout、学生 readiness 或 TA feedback 不清楚时，先修材料。
\end{enumerate}

最小形式可以写成：

\begin{examplebox}[一个最小 trial-teaching routing 逻辑]
\begin{lstlisting}[style=pythonstyle]
if privacy_boundary_status != "clear":
    blocked_before_trial
elif notebook_runtime_status == "fails":
    blocked_before_trial
elif rubric_calibration_status != "calibrated":
    calibrate_rubric
elif handout_clarity_status != "clear":
    revise_materials
elif student_readiness_status == "low":
    revise_materials
elif ta_feedback_status != "ready":
    revise_materials
elif revision_load_score > 0.45:
    revise_materials
else:
    ready_for_rollout
\end{lstlisting}
\end{examplebox}

这个顺序很重要。隐私边界和 notebook 运行失败是硬阻塞；rubric calibration 是评分一致性的前提；handout 和 feedback 格式则决定学生是否能真的照着材料完成项目。

\section{学生 handout 应该包含什么}

把 Capstone 正式发给学生时，推荐至少给出一份 4 到 6 页 handout，包含：

\begin{enumerate}
    \item 项目目标：学生要回答什么科学问题，而不是只训练一个模型；
    \item 时间线：proposal、pilot、project board、draft report、final delivery 和 presentation 的日期；
    \item 数据边界：哪些数据可以进入仓库，哪些不能上传到外部服务；
    \item notebook 要求：必须能从头运行，必须说明环境和随机种子；
    \item rubric：六个评分维度、评分前 gate 和最低线；
    \item AI 使用规范：鼓励、限制、禁止和必须披露的场景；
    \item 求助路径：什么时候问助教，什么时候必须 human signoff；
    \item 提交清单：报告、notebook、图表、AI-use statement、Git 记录和展示材料。
\end{enumerate}

这份 handout 的目标不是把所有可能性写完，而是让学生在遇到不确定时知道下一步该往哪里走。

\section{助教评分校准也需要材料}

如果课程有助教，rubric 还需要一份 TA calibration package。它至少应该包含：

\begin{itemize}
    \item 两到三个匿名样例项目包；
    \item 每个样例的建议 route 和关键证据；
    \item 常见扣分点和不可扣分点；
    \item AI-use statement 的合格与不合格样例；
    \item 对同一个项目的 2 到 3 条示范 feedback；
    \item 当助教无法判断时的升级路径。
\end{itemize}

没有这一步，rubric 很容易从“统一标准”变成“每个评分者心里各有一套标准”。

\section{配套 notebook 做了什么}

本章 notebook 采用的是一个 teaching-first 的最小设计：

\begin{itemize}
    \item 先读取 trial-teaching feedback 表，并统计 student readiness、runtime、rubric calibration 和 privacy boundary；
    \item 用 \texttt{revision\_load\_score} 做一个 naive rollout baseline；
    \item 再实现一个带 privacy gate、runtime gate、rubric gate 和 materials gate 的 trial workflow；
    \item 输出 ready、revise materials、calibrate rubric 和 blocked 四个队列；
    \item 为每个非 ready 环节生成课程修订 checklist；
    \item 最后生成一份 trial-teaching assistant prompt 和 rollout snapshot。
\end{itemize}

这份 notebook 不是在把课程管理自动化，而是在提醒我们：课程材料也需要像科研代码一样测试、记录和迭代。

\section{AI 助手如何帮助你}

在这个案例里，AI 助手最适合帮助你：

\begin{itemize}
    \item 把 Capstone handout 改写成更清楚的学生任务单；
    \item 根据试教反馈归类问题：材料问题、评分问题、运行问题还是边界问题；
    \item 生成助教 norming session 的样例反馈；
    \item 把学生反馈整理成下一轮课程修订任务；
    \item 检查 AI 使用规范里是否有模糊或互相矛盾的句子。
\end{itemize}

但 AI 助手不应该替课程团队决定隐私边界、评分政策或学生最终成绩。它可以帮助整理证据和发现不一致，最终课程规则仍然需要由人类教师确认。

\section{练习}

\begin{exercisebox}[基础练习]
把教学数据中一个 \texttt{ready\_for\_rollout} 环节的 \texttt{notebook\_runtime\_status} 改成 \texttt{fails}。重新运行 notebook，观察它为什么会直接进入试教前阻塞队列。
\end{exercisebox}

\begin{exercisebox}[进阶练习]
新增一条试教记录：学生 readiness 是 \texttt{ready}，handout 清楚，notebook 可以运行，但 \texttt{rubric\_calibration\_status} 是 \texttt{uncalibrated}。预测它会落到哪个 route，并说明为什么评分校准应该早于正式 rollout。
\end{exercisebox}

\begin{exercisebox}[开放练习]
为你自己的课程项目写一份一页试教计划。至少包含：参与学生、运行环境、要收集的反馈、助教校准方式、隐私边界和下一轮修订负责人。
\end{exercisebox}

\section{本章小结}

Capstone 材料写完以后，还需要经过课堂运行检验。本章把学生 readiness、handout clarity、notebook runtime、rubric calibration、TA feedback、privacy boundary 和 revision load 放进同一张 trial-teaching feedback card。

如果课程团队愿意在正式开课前做这一次小型系统调试，学生拿到的就不只是一个期末项目题目，而是一条更清楚、更安全、也更可持续的科研入门路径。
